\documentclass[10pt,a4paper]{amsart}
\usepackage[margin=19mm]{geometry}
\usepackage{amsmath,amssymb,mathtools}
\usepackage[scr=rsfs,cal=euler]{mathalfa}
\usepackage{xfrac}
\usepackage[unicode,hidelinks,bookmarks=false]{hyperref}
\newcommand{\ie}{i.e.\ }
\newcommand{\cf}{cf.\ }
\newcommand{\resp}{respectively\ }
\newcommand{\Subsection}{\subsection}
\providecommand{\nobreakdash}{\nobreak\hbox{-}}
\setlength{\emergencystretch}{2em}
\multlinegap0pt
\usepackage{amssymb}
\usepackage{booktabs}
\usepackage{enumerate}
\usepackage{mathtools}
\usepackage{dsfont}
\usepackage{xcolor}
\newtheorem{lem}{Lemma}
\newcommand{\N}{\mathds{N}}
\newcommand{\de}{\partial}
\newcommand{\pk }{para-K\"{a}hler }
\title{Toric para-K\"{a}hler-Einstein manifolds immersed in para-K\"{a}hler space forms}
\author{Gianni Manno, Filippo Salis}
\date{Source: arXiv:2510.19625v1 (CC BY 4.0)}
\begin{document}
\maketitle

\noindent\emph{Source and licence.} Toric para-K\"{a}hler-Einstein manifolds immersed in para-K\"{a}hler space forms. Version arXiv:2510.19625v1 (CC BY 4.0), distributed by arXiv under the Creative Commons Attribution 4.0 International licence, http://creativecommons.org/licenses/by/4.0/, as recorded in the arXiv licence metadata for this exact version and shown on the abstract page https://arxiv.org/abs/2510.19625v1. Full licence notice: \texttt{licenses/arxiv-2510.19625-CC-BY-4.0.txt}.\par\smallskip

\noindent\textit{Editorial scope.} Counted unit: the article's lem lem.polyn (source0.tex lines 597--607). The article's complete original proof of it is delivered with it (source0.tex lines 608--684, one \texttt{proof} environment of 77 lines). No mathematics is written, repaired or re-derived: the statement and the proof are the article's own lines, in the article's own notation, and no retained line is substituted or re-flowed.  No further source-local context is needed: every cross-reference of every retained line resolves inside the two delivered spans.  Nothing was omitted from any retained span, so the package carries no omission record.  No retained line contains a citation, so the wrapper carries no bibliography and no bibliography entry.  No retained line contains a figure environment, an image-inclusion command or drawing code, so no image is delivered and the package declares no asset dependency.  Editorial formatting only, in addition: an AMS \texttt{amsart} preamble that restates the article's own declarations and macros used by the retained lines, namely the environments \texttt{lem} on separate counters, together with the standard packages those lines need.  The retained environments print in this document as Lemma 1 in order of appearance, whereas the article prints the same items with the numbering of its own full document.  The complete native source of the article as distributed in the arXiv e-print for this exact version is bundled unchanged as \texttt{sources/187586/source0.tex}.
\begin{lem}\label{lem.polyn}
Let $\phi$ be a toric-invariant \pk potential.
If $u_i$ and $v_i$ are some analytic functions such that
\begin{equation}\label{eq:1}
\sum_{i=1}^N u_i(\xi)v_i(\eta)=\phi(\xi_1\eta_1,\mathellipsis,\xi_n\eta_n),\end{equation}
where  $N\in\N$ is the smallest possible,
then they are polynomials. Hence, $\phi$ is a polynomial in the variables 
\begin{equation}\label{eq.x.xi.eta}
x_i=\xi_i\eta_i\,.
\end{equation}
\end{lem}

\begin{proof}
By differentiating \eqref{eq:1} and evaluating at $\xi=0$, we obtain the following (compatible) linear system
\begin{equation}\label{system}
\begin{pmatrix}
\frac{\de^{|I_1|} u_1}{\de\xi^{I_1}}(0) &\dots&\frac{\de^{|I_1|}  u_N}{\de\xi^{I_1}}(0) \\
\vdots& &\vdots\\
\frac{\de^{|I_N|}  u_1}{\de\xi^{I_N}}(0) &\dots&\frac{\de^{|I_N|}  u_N}{\de\xi^{I_N}}(0) \\
\end{pmatrix}
\begin{pmatrix}
v_1(\eta)\\
\vdots\\
v_N(\eta)
\end{pmatrix}
=\begin{pmatrix}
\eta^{I_1}\frac{\de^{|I_1|}  \phi}{\de x^{I_1}}(0)\\
\vdots\\
\eta^{I_N}\frac{\de^{|I_N|}  \phi}{\de x^{I_N}}(0)
\end{pmatrix},
\end{equation}
where $I_i\in\N^n$ are arbitrary chosen multi-indices. As a first step, we notice that there exist some $I_1,\mathellipsis,I_N\in\N^n$ such that
$$\mathrm{rank} \begin{pmatrix}
\frac{\de^{|I_1|} u_1}{\de\xi^{I_1}}(0) &\dots&\frac{\de^{|I_1|} u_N}{\de\xi^{I_1}}(0) \\
\vdots& &\vdots\\
\frac{\de^{|I_N|} u_1}{\de\xi^{I_N}}(0) &\dots&\frac{\de^{|I_N|} u_N}{\de\xi^{I_N}}(0) \\
\end{pmatrix}>0.
$$ 
Indeed, if, on the contrary, the previous matrix had rank equal to 0 for any $I_1,\mathellipsis,I_N$, then, in view of the analyticity of $u_1,\dots,u_N$, we would have
$$u_1=\mathellipsis=u_N\equiv 0$$  and $\phi\equiv0$ would  not be a potential of a symplectic form.

\smallskip\noindent
Now, if
$$\mathrm{rank} \begin{pmatrix}
\frac{\de^{|I_1|}  u_1}{\de\xi^{I_1}}(0) &\dots&\frac{\de^{|I_1|}  u_N}{\de\xi^{I_1}}(0) \\
\vdots& &\vdots\\
\frac{\de^{|I_N|}  u_1}{\de\xi^{I_N}}(0) &\dots&\frac{\de^{|I_N|}  u_N}{\de\xi^{I_N}}(0) \\
\end{pmatrix}=N$$
for some $I_1,\mathellipsis,I_N\in\N^n$, then, being $(v_1,\mathellipsis,v_N)$ a solution of the linear system \eqref{system}, $v_1,\mathellipsis,v_N$ are polynomials.


\noindent
To conclude, we take into account the case in which
$$\mathrm{rank} \begin{pmatrix}
\frac{\de^{|I_1|} u_1}{\de\xi^{I_1}}(0) &\dots&\frac{\de^{|I_1|} u_N}{\de\xi^{I_1}}(0) \\
\vdots& &\vdots\\
\frac{\de^{|I_N|} u_1}{\de\xi^{I_N}}(0) &\dots&\frac{\de^{|I_N|} u_N}{\de\xi^{I_N}}(0) \\
\end{pmatrix}<N$$
for any $I_1,\mathellipsis,I_N\in\N^n$. In view of what we said above, we can assume that
\begin{equation}\label{rank1}
\mathrm{rank} \begin{pmatrix}
\frac{\de^{|I_1|} u_1}{\de\xi^{I_1}}(0) &\dots&\frac{\de^{|I_1|} u_\rho}{\de\xi^{I_1}}(0) \\
\vdots& &\vdots\\
\frac{\de^{|I_\rho|} u_1}{\de\xi^{I_\rho}}(0) &\dots&\frac{\de^{|I_\rho|} u_\rho}{\de\xi^{I_\rho}}(0) \\
\end{pmatrix}=\rho>0\end{equation}
for some suitable  $I_1,\mathellipsis,I_\rho\in\N^n$ and
\begin{equation}\label{rank2}
\mathrm{rank} \begin{pmatrix}
\frac{\de^{|I_1|} u_1}{\de\xi^{I_1}}(0) &\dots&\frac{\de^{|I_1|}  u_{\rho+1}}{\de\xi^{I_1}}(0) \\
\vdots& &\vdots\\
\frac{\de^{|I_\rho|}  u_1}{\de\xi^{I_\rho}}(0) &\dots&\frac{\de^{|I_\rho|} u_{\rho+1}}{\de\xi^{I_\rho}}(0) \\
\frac{\de^{|J|} u_1}{\de\xi^{J}}(0) &\dots&\frac{\de^{|J|} u_{\rho+1}}{\de\xi^{J}}(0) \\
\end{pmatrix}=\rho\end{equation}
for any $J\in\N^n$.
It follows from the Laplace's expansion w.r.t. the last row of the matrix in \eqref{rank2}, and by considering also \eqref{rank1}, that
\begin{equation*}%\label{taylor}
\frac{\de^{|J|} u_{\rho+1}}{\de\xi^{J}}(0)=-\frac{K_1}{K_{\rho+1}} \frac{\de^{|J|} u_1}{\de\xi^{J}}(0)-\mathellipsis -\frac{K_\rho}{K_{\rho+1}} \frac{\de^{|J|} u_\rho}{\de\xi^{J}}(0),\qquad \forall J\in\N^n,
\end{equation*}
where $K_i$ denotes the algebraic complement of $i$-th element of the last row of the matrix in \eqref{rank2}.
Since
$$
u_{\rho+1}(\xi)=\sum_{J\in\N^n}\frac{\de^{|J|} u_{\rho+1}}{\de\xi^{J}}(0)\  \frac{\xi^J}{J!} =-\frac{K_1}{K_{\rho+1}} u_1(\xi)-\mathellipsis-\frac{K_\rho}{K_{\rho+1}} u_\rho(\xi)\,,
$$
by taking suitable linear combinations $\tilde v_1,\mathellipsis,\tilde v_\rho$ of $v_1,\mathellipsis,v_n$, we have 
$$\sum_{i=1}^\rho u_i\tilde v_i=\phi.$$ Hence, $N$ in \eqref{eq:1} would not be the smallest integer, as assumed.

\smallskip\noindent
The same reasoning can be applied to the functions $u_i$.
\end{proof}

\end{document}
